\chapter{The quantum Fourier transform}
\label{ch:qft}

This chapter writes the quantum Fourier transform on $n$ qubits as a unitary circuit
of one- and two-qubit gates, in the form of \cite[Fig.~5.1]{NielsenChuang2010}, and
proves that its Kraus operator is the Fourier matrix. It then schedules the same
gates in $2n-1$ layers of disjoint gates, followed by one layer of swaps, which gives
a circuit of depth at most $2n$ for the same matrix. The chapter uses the circuits
of Chapter~\ref{ch:circuits} on the geometry of sites and pairs.

\section{Qubits and gates}
\label{sec:qft_qubits}

\begin{definition}[Qubits and the binary encoding]
    \label{def:qft_qubits}
    \lean{QuantumCircuit.Qubit.Q, QuantumCircuit.Qubit.encode,
        QuantumCircuit.Qubit.encode_apply, QuantumCircuit.Qubit.geometry}
    \leanok
    \uses{def:sites_config, def:circuits_local_gate}
    The system of $n$ qubits has the sites $0,\dots,n-1$, each with basis
    $\{0,1\}$. A configuration $x$ encodes the number
    $\langle x\rangle=\sum_ix_i2^{n-1-i}$, the site $0$ most significant; this is a
    bijection onto $\{0,\dots,2^n-1\}$. The circuits of this chapter use the
    geometry of sites and pairs on the $n$ qubits, every cell of cost $1$.
\end{definition}

\begin{definition}[Gates]
    \label{def:qft_gates}
    \lean{QuantumCircuit.Qubit.hadamard, QuantumCircuit.Qubit.phaseShift,
        QuantumCircuit.Qubit.controlled, QuantumCircuit.Qubit.swap,
        QuantumCircuit.Qubit.siteLayer, QuantumCircuit.Qubit.pairLayer}
    \leanok
    \uses{def:qft_qubits, lem:circuits_layer_single}
    The Hadamard gate is $H_{yx}=(-1)^{xy}/\sqrt2$, the phase gate is
    $R_k=\operatorname{diag}(1,e^{2\pi i/2^k})$, the controlled gate of a matrix $U$
    is $\lvert0\rangle\langle0\rvert\otimes1+\lvert1\rangle\langle1\rvert\otimes U$,
    the control first, and the swap exchanges two qubits. A gate on one qubit or on
    an ordered pair of distinct qubits gives the layer with this single gate.
\end{definition}

\begin{lemma}[The gates are unitary]
    \label{lem:qft_gates_unitary}
    \lean{QuantumCircuit.Qubit.hadamard_mem_unitaryGroup,
        QuantumCircuit.Qubit.phaseShift_mem_unitaryGroup,
        QuantumCircuit.Qubit.controlled_mem_unitaryGroup,
        QuantumCircuit.Qubit.swap_mem_unitaryGroup,
        QuantumCircuit.Qubit.forall_isUnitary_siteLayer,
        QuantumCircuit.Qubit.forall_isUnitary_pairLayer,
        QuantumCircuit.Qubit.kraus_siteLayer, QuantumCircuit.Qubit.kraus_pairLayer,
        QuantumCircuit.Qubit.nonempty_index_siteLayer,
        QuantumCircuit.Qubit.nonempty_index_pairLayer}
    \leanok
    \uses{def:qft_gates, lem:circuits_local_gate_unitary}
    The four gates are unitary, so the layers of one of them are locally unitary.
    The single Kraus operator of such a layer is the placement of the gate.
\end{lemma}
\begin{proof}\leanok
    \uses{lem:circuits_layer_single}
    Direct computation for the four matrices; the controlled gate of a unitary
    matrix is block diagonal with unitary blocks. The second statement is
    Lemma~\ref{lem:circuits_layer_single}.
\end{proof}

\begin{lemma}[Gates on basis vectors]
    \label{lem:qft_gates_vectors}
    \lean{QuantumCircuit.Qubit.cut_update, QuantumCircuit.Qubit.place_site_mulVec,
        QuantumCircuit.Qubit.place_diagonal_mulVec,
        QuantumCircuit.Qubit.controlled_phaseShift}
    \leanok
    \uses{def:qft_gates, def:sites_place}
    A gate $K$ placed on one qubit $a$ maps a vector $\psi$ to the vector whose value
    at $\sigma$ is $\sum_bK_{\sigma_a,b}\,\psi(\sigma[a\mapsto b])$. A diagonal gate
    $\operatorname{diag}(d)$ on a region multiplies the value at $\sigma$ by
    $d(\sigma|_D)$. The controlled phase gate of $R_k$ is the diagonal gate
    $(c,t)\mapsto e^{2\pi i\,ct/2^k}$.
\end{lemma}
\begin{proof}\leanok
    The formula for a placed matrix acting on a vector, from
    Definition~\ref{def:sites_place}, with the configurations of $\{a\}$ identified
    with $\{0,1\}$.
\end{proof}

\section{The Fourier transform}
\label{sec:qft_circuit}

\begin{definition}[The Fourier matrix and its circuit]
    \label{def:qft_circuit}
    \lean{QuantumCircuit.Qubit.dft, QuantumCircuit.Qubit.rotationLayer,
        QuantumCircuit.Qubit.qubitLayers, QuantumCircuit.Qubit.coreLayers,
        QuantumCircuit.Qubit.rev_ne_of_lt_half, QuantumCircuit.Qubit.swapLayers,
        QuantumCircuit.Qubit.qftCore, QuantumCircuit.Qubit.qft}
    \leanok
    \uses{def:qft_gates, def:circuits_circuit}
    The Fourier matrix on $n$ qubits is
    \begin{align}
        F_{yx}&=2^{-n/2}\,e^{2\pi i\langle y\rangle\langle x\rangle/2^n}
        \notag
    \end{align}
    \cite[eq.~(5.2)]{NielsenChuang2010}. For each qubit $i=0,\dots,n-1$ in turn, the
    circuit applies the Hadamard gate to $i$ and then, for $j=i+1,\dots,n-1$, the
    controlled gate of $R_{j-i+1}$ with control $j$ and target $i$. Its Fourier
    circuit follows these gates by the swaps of the qubits $k$ and $n-1-k$ for
    $k<\lfloor n/2\rfloor$. Every layer holds one gate.
\end{definition}

\begin{theorem}[The circuit computes the Fourier transform]
    \label{thm:qft_kraus}
    \lean{QuantumCircuit.Qubit.kraus_qftCore, QuantumCircuit.Qubit.kraus_qft,
        QuantumCircuit.Qubit.U1, QuantumCircuit.Qubit.U1_siteLayer,
        QuantumCircuit.Qubit.U1_pairLayer, QuantumCircuit.Qubit.prod_mulVec,
        QuantumCircuit.Qubit.rotation_mulVec, QuantumCircuit.Qubit.swap_mulVec,
        QuantumCircuit.Qubit.kraus_eq_U1_of_mem, QuantumCircuit.Qubit.phase,
        QuantumCircuit.Qubit.qstate, QuantumCircuit.Qubit.qstate',
        QuantumCircuit.Qubit.neg_one_pow_eq_exp, QuantumCircuit.Qubit.hadamard_qstate,
        QuantumCircuit.Qubit.rotation_qstate', QuantumCircuit.Qubit.foldl_rotations,
        QuantumCircuit.Qubit.sum_filter_finRange, QuantumCircuit.Qubit.phase_eq,
        QuantumCircuit.Qubit.foldl_qubitLayers, QuantumCircuit.Qubit.qstate_zero,
        QuantumCircuit.Qubit.foldl_coreLayers, QuantumCircuit.Qubit.swapped,
        QuantumCircuit.Qubit.foldl_swapLayers, QuantumCircuit.Qubit.encode_rev_cast,
        QuantumCircuit.Qubit.encode_cast, QuantumCircuit.Qubit.pow_ratio,
        QuantumCircuit.Qubit.exponent_eq, QuantumCircuit.Qubit.qstate_n}
    \leanok
    \uses{def:qft_circuit, lem:circuits_seq}
    Along its single branch the circuit without the swaps has the Kraus operator
    whose entry at $(y,x)$ is $F_{\bar y x}$, where $\bar y_i=y_{n-1-i}$. The Kraus
    operator of the Fourier circuit is $F$.
\end{theorem}
\begin{proof}\leanok
    \uses{lem:qft_gates_unitary, lem:qft_gates_vectors}
    By Lemma~\ref{lem:circuits_seq} the Kraus operator is the product of the placed
    gates, and it suffices to follow the column $e_x$ of every basis vector. After
    the gates for the qubits $0,\dots,i-1$ the vector is the product state whose
    factor at each processed qubit $k$ is
    $2^{-1/2}(\lvert0\rangle+e^{2\pi i\,\phi_k}\lvert1\rangle)$ with
    $\phi_k=\sum_{j\geq k}x_j2^{-(j-k+1)}$, the remaining qubits still in their
    basis states. The Hadamard gate on $i$ produces the phase $x_i/2$, and the
    controlled gate of $R_{j-i+1}$ adds $x_j2^{-(j-i+1)}$, since the control $j$ is
    still in its basis state. At the end the entry at $y$ is
    $2^{-n/2}\prod_ke^{2\pi i\,y_k\phi_k}$, and
    $\sum_ky_k\phi_k$ differs from $\langle\bar y\rangle\langle x\rangle/2^n$ by an
    integer: the terms $x_jy_k2^{j-k-1}$ with $j<k$ omitted from the product are
    integers. The swaps reverse the order of the qubits.
\end{proof}

\begin{theorem}[Unitarity and cost of the Fourier circuit]
    \label{thm:qft_cost}
    \lean{QuantumCircuit.Qubit.isUnbranched_qft, QuantumCircuit.Qubit.forall_isUnitary_qft,
        QuantumCircuit.Qubit.nonempty_index_and_gateCount_of_mem,
        QuantumCircuit.Qubit.gateCount_seq_eq_length, QuantumCircuit.Qubit.length_qubitLayers,
        QuantumCircuit.Qubit.length_coreLayers, QuantumCircuit.Qubit.gateCount_qftCore,
        QuantumCircuit.Qubit.gateCount_qft, QuantumCircuit.Qubit.depth_qft_le}
    \leanok
    \uses{def:qft_circuit, def:circuits_costs, def:circuits_local}
    The Fourier circuit is unbranched and locally unitary. Without the swaps it uses
    $n(n+1)/2$ gates, and with them $n(n+1)/2+\lfloor n/2\rfloor$, which also bounds
    its depth.
\end{theorem}
\begin{proof}\leanok
    \uses{lem:circuits_seq, lem:qft_gates_unitary, thm:circuits_cost_subadditive}
    Every layer holds one gate of cost $1$ and has one outcome, so by
    Lemma~\ref{lem:circuits_seq} the gate count is the number of layers: $n-i$ for
    the qubit $i$, in total $n(n+1)/2$, and $\lfloor n/2\rfloor$ swaps. The depth is
    at most the gate count.
\end{proof}

\section{The Fourier transform in linear depth}
\label{sec:qft_parallel}

\begin{definition}[The scheduled Fourier circuit]
    \label{def:qft_parallel}
    \lean{QuantumCircuit.Qubit.qftV, QuantumCircuit.Qubit.qftGate,
        QuantumCircuit.Qubit.timeCells, QuantumCircuit.Qubit.mem_cell_inl,
        QuantumCircuit.Qubit.mem_cell_inr, QuantumCircuit.Qubit.timeCells_pairwiseDisjoint,
        QuantumCircuit.Qubit.timeLayer, QuantumCircuit.Qubit.swapCell,
        QuantumCircuit.Qubit.swapCells, QuantumCircuit.Qubit.swapV,
        QuantumCircuit.Qubit.swapGate, QuantumCircuit.Qubit.swapCells_pairwiseDisjoint,
        QuantumCircuit.Qubit.swapAll, QuantumCircuit.Qubit.qftPar}
    \leanok
    \uses{def:qft_circuit, def:circuits_of_cells}
    Place the Hadamard gate on the qubit $i$ at time $2i$ and the controlled gate of
    $R_{j-i+1}$ with control $j$ and target $i<j$ at time $i+j$. The cells used at
    one time $t$ are pairwise disjoint: the pairs $\{i,j\}$ with $i+j=t$ are
    disjoint, and the qubit $t/2$ lies in none of them. The scheduled circuit is the
    sequence of the layers of the gates at the times $0,\dots,2n-2$, followed by the
    layer of all swaps of $k$ and $n-1-k$ for $k<\lfloor n/2\rfloor$.
\end{definition}

\begin{theorem}[The scheduled circuit computes the Fourier transform in depth $2n$]
    \label{thm:qft_parallel}
    \lean{QuantumCircuit.Qubit.kraus_qftPar, QuantumCircuit.Qubit.depth_qftPar_le,
        QuantumCircuit.Qubit.pstate, QuantumCircuit.Qubit.pstate_congr,
        QuantumCircuit.Qubit.hadamard_pstate, QuantumCircuit.Qubit.rotation_pstate,
        QuantumCircuit.Qubit.xval, QuantumCircuit.Qubit.doneSet, QuantumCircuit.Qubit.θt,
        QuantumCircuit.Qubit.tstate, QuantumCircuit.Qubit.tstate_zero,
        QuantumCircuit.Qubit.rotValid, QuantumCircuit.Qubit.rot_ne,
        QuantumCircuit.Qubit.rotCell, QuantumCircuit.Qubit.rList,
        QuantumCircuit.Qubit.contrib, QuantumCircuit.Qubit.contrib_nil,
        QuantumCircuit.Qubit.contrib_cons, QuantumCircuit.Qubit.rList_mulVec,
        QuantumCircuit.Qubit.hList, QuantumCircuit.Qubit.timeList,
        QuantumCircuit.Qubit.mem_rList, QuantumCircuit.Qubit.timeList_nodup,
        QuantumCircuit.Qubit.timeList_toFinset, QuantumCircuit.Qubit.θt_succ,
        QuantumCircuit.Qubit.contrib_finRange, QuantumCircuit.Qubit.timeList_mulVec,
        QuantumCircuit.Qubit.kraus_timeLayer, QuantumCircuit.Qubit.nonempty_index_timeLayer,
        QuantumCircuit.Qubit.U1_timeLayer, QuantumCircuit.Qubit.swapCell_injective,
        QuantumCircuit.Qubit.kraus_swapAll, QuantumCircuit.Qubit.nonempty_index_swapAll,
        QuantumCircuit.Qubit.U1_swapAll, QuantumCircuit.Qubit.list_prod_mulVec,
        QuantumCircuit.Qubit.foldl_timeLayers, QuantumCircuit.Qubit.tstate_final,
        QuantumCircuit.Qubit.kraus_eq_U1_of_mem_par}
    \leanok
    \uses{def:qft_parallel, def:circuits_costs}
    The Kraus operator of the scheduled circuit is the Fourier matrix $F$, and its
    depth is at most $2n$.
\end{theorem}
\begin{proof}\leanok
    \uses{thm:circuits_of_cells, lem:circuits_seq, lem:qft_gates_vectors,
        thm:qft_kraus, thm:circuits_cost_subadditive}
    By Theorem~\ref{thm:circuits_of_cells} the Kraus operator of a layer is the
    product of the placements of its gates, which commute since their cells are
    disjoint. Follow the column $e_x$ as in Theorem~\ref{thm:qft_kraus}. After the
    times $0,\dots,t-1$ the vector is the product state in which the qubits $k$ with
    $2k<t$ carry the factor
    $2^{-1/2}(\lvert0\rangle+e^{2\pi i\,\phi_k^{(t)}}\lvert1\rangle)$, where
    $\phi_k^{(t)}=x_k/2+\sum_{k<j,\,k+j<t}x_j2^{-(j-k+1)}$, and the other qubits are
    in their basis states. At time $t$ the Hadamard gate acts on a qubit still in its
    basis state, and each controlled gate with control $j$ and target $i$, $i+j=t$,
    acts while $j$ is in its basis state, since $2j>t$. At time $2n-1$ the phases are
    those of Theorem~\ref{thm:qft_kraus}, and the swaps reverse the order of the
    qubits. Each of the $2n$ layers has depth at most $1$, and the depth of a
    sequence is at most the sum of the depths of its layers.
\end{proof}
